\documentclass[12pt,a4paper]{article}

\PassOptionsToPackage{hyphens}{url}   % let long URLs break at / and -

\usepackage[a4paper,margin=1in]{geometry}
\usepackage{graphicx}
\usepackage{booktabs}
\usepackage{amsmath,amssymb}
\usepackage{float}
\usepackage{hyperref}
\usepackage{xcolor}
\usepackage{listings}
\usepackage{enumitem}
\usepackage{fancyhdr}
\usepackage{longtable}
\usepackage{array}

% Scale a table down to the text width only if it would otherwise overflow.
\newsavebox{\fitbox}
\newenvironment{fitwidth}
  {\begin{lrbox}{\fitbox}}
  {\end{lrbox}\ifdim\wd\fitbox>\linewidth\resizebox{\linewidth}{!}{\usebox{\fitbox}}\else\usebox{\fitbox}\fi}
% Let long code identifiers (e.g. load\_breast\_cancer) wrap at underscores instead of running into the margin.
\DeclareRobustCommand{\_}{\textunderscore\allowbreak}
\setlength{\emergencystretch}{3em}

\hypersetup{colorlinks=true,linkcolor=blue,urlcolor=blue}

\pagestyle{fancy}
\fancyhf{}
\lhead{ICS1512}
\rhead{Machine Learning Algorithms Laboratory}
\cfoot{\thepage}

\lstset{
language=Python,
basicstyle=\ttfamily\small,
numbers=left,
frame=single,
breaklines=true,
showstringspaces=false,
keywordstyle=\color{blue},
commentstyle=\color{gray},
stringstyle=\color{purple}
}

\begin{document}

\noindent\begin{tabular}{|p{4cm}|p{\dimexpr\linewidth-4cm-4\tabcolsep-3\arrayrulewidth\relax}|}
\hline
Experiment No. & 8 \\ \hline
Experiment Title & Clustering Human Activity Recognition Data using K-Means, DBSCAN and Hierarchical Clustering
\footnote{\textbf{GitHub Repository: }\url{https://github.com/Nightingales21/ICS1512-Machine-Learning-Labratory}}\\ \hline
Student Name & Nitin \\ \hline
Register Number & 312224700142 \\ \hline
Faculty & Dr. Poreddy Ajay Kumar Reddy \\ \hline
Submission Date & \today \\ \hline
\end{tabular}

\section{Aim and Objective}
To implement and analyse three clustering algorithms on the Human Activity
Recognition (HAR) dataset:
\begin{itemize}[nosep]
\item \textbf{Model A:} K-Means clustering, with $k$ selected by the Elbow
method and silhouette analysis;
\item \textbf{Model B:} DBSCAN (Density-Based Spatial Clustering of
Applications with Noise), with $\epsilon$ and \textit{minPts} tuned;
\item \textbf{Model C:} Hierarchical Agglomerative Clustering (HAC), with a
dendrogram and a comparison of linkage criteria.
\end{itemize}
The clusters produced by each algorithm are visualised in 2-D, compared
against the ground-truth activity labels (which the algorithms never see),
and evaluated with both internal and external clustering metrics.

\section{Dataset Description}
\begin{longtable}{|p{6cm}|p{8cm}|}
\hline
Dataset Name & Human Activity Recognition Using Smartphones (UCI HAR) \\ \hline
Dataset Source & UCI Machine Learning Repository (dataset 240), retrieved via
OpenML (\texttt{data\_id=1478}) and stored as \texttt{har\_dataset.zip} \\ \hline
Number of Samples & 10{,}299 windows (train and test partitions pooled, since
clustering is unsupervised) \\ \hline
Number of Features & 561 engineered time- and frequency-domain features \\ \hline
Number of Classes & 6 activities (used only for external validation) \\ \hline
Subjects & 30 volunteers aged 19--48 \\ \hline
Sensors & 3-axis accelerometer and gyroscope at 50\,Hz, windowed at 2.56\,s
(128 readings/window) with 50\% overlap \\ \hline
Missing Values & 0 (verified programmatically) \\ \hline
\end{longtable}

\noindent Class balance: LAYING 1{,}944, STANDING 1{,}906, SITTING 1{,}777,
WALKING 1{,}722, WALKING\_UPSTAIRS 1{,}544, WALKING\_DOWNSTAIRS 1{,}406 --
mildly imbalanced (a 1.38:1 ratio between the largest and smallest activity),
which matters for clustering because algorithms that favour equally sized
clusters are at a slight disadvantage.

\section{Preprocessing Steps}
\begin{itemize}[nosep]
\item \textbf{Missing value handling:} none required; the published feature
matrix has no missing entries.
\item \textbf{Categorical encoding:} none required for the features (all 561
are continuous). Activity names are mapped to integer indices purely so that
the external metrics (ARI, NMI) can be computed.
\item \textbf{Standardization:} the published features are already bounded to
$[-1, 1]$, but their spreads differ, so every feature was standardized to zero
mean and unit variance with \texttt{StandardScaler}. All three algorithms here
are distance-based, so an unstandardized feature with a wider spread would
silently dominate every distance computation.
\item \textbf{Dimensionality reduction:} PCA is used for two distinct
purposes -- 2-D projections for the figures, and a 10-component space in
which DBSCAN is tuned (Section~\ref{sec:dbscan} explains why DBSCAN needs
this). K-Means
and HAC are run on the full 561-dimensional standardized matrix.
\end{itemize}

\noindent PCA needs \textbf{104 of 561 components} to retain 95\% of the
variance, but the first component alone already explains \textbf{50.7\%} --
a strong sign that one dominant axis (body motion intensity) separates most
of the data.

\section{Exploratory Data Analysis}

\begin{figure}[H]
\centering
\includegraphics[width=\linewidth]{figures/eda_har.png}
\caption{EDA summary: class distribution, PCA scree plot, PCA and t-SNE
projections coloured by true activity, a per-activity feature distribution,
and the correlation structure of the first 60 features.}
\label{fig:eda}
\end{figure}
\textbf{Inference:} The PCA projection (panel 3) shows the single most
important structural fact about this dataset: the six activities form
\emph{two} widely separated super-clusters -- the three static postures
(LAYING, SITTING, STANDING) and the three dynamic gaits (WALKING,
WALKING\_UPSTAIRS, WALKING\_DOWNSTAIRS) -- with a large empty gap between
them. Within each super-cluster the individual activities overlap heavily:
SITTING and STANDING are nearly superimposed, and the three walking variants
form one elongated band. The t-SNE view (panel 4) separates LAYING cleanly and
teases the walking variants apart a little, but still leaves SITTING and
STANDING interleaved. This two-level structure predicts the central result of
the whole experiment: unsupervised methods will find the 2-cluster split
easily and the 6-activity split only partially.

\section{Experimental Setup}
\begin{center}
\begin{minipage}{\linewidth}\centering
\begin{fitwidth}
\begin{tabular}{ll}
\toprule
Python Version & 3.11.15 \\
Scikit-Learn Version & 1.9.0 \\
NumPy / Pandas Version & 2.4.6 / 3.0.5 \\
SciPy Version & 1.17.1 \\
IDE & Jupyter Notebook / VS Code \\
Operating System & Windows 11 \\
Hardware & Standard lab workstation (CPU only) \\
\bottomrule
\end{tabular}
\end{fitwidth}
\end{minipage}
\end{center}

\vspace{0.3cm}
\textbf{Environment activation command (as mandated by the lab manual):}
\begin{lstlisting}[language=bash]
source /opt/anaconda3/bin/activate anaconda-navigation
\end{lstlisting}

\section{Model A: K-Means and the Elbow Method}

\begin{center}
\begin{minipage}{\linewidth}\centering
\textbf{Table 1: K-Means Elbow Method Results (as required by the lab manual)}\par\smallskip
\begin{fitwidth}
\begin{tabular}{lcccccc}
\toprule
$k$ & WCSS (Inertia) & Silhouette & Davies--Bouldin & Calinski--Harabasz & ARI & NMI \\
\midrule
2 & 3{,}272{,}857 & \textbf{0.3946} & \textbf{1.071} & \textbf{7880.8} & 0.3296 & 0.5455 \\
3 & 2{,}921{,}078 & 0.3137 & 1.787 & 5034.5 & 0.3317 & 0.5183 \\
4 & 2{,}781{,}602 & 0.1524 & 2.342 & 3696.3 & 0.2990 & 0.4694 \\
5 & 2{,}654{,}790 & 0.1321 & 2.411 & 3027.3 & 0.2912 & 0.4607 \\
6 & 2{,}577{,}180 & 0.1095 & 2.384 & 2556.5 & 0.4196 & \textbf{0.5593} \\
7 & 2{,}519{,}774 & 0.0808 & 2.675 & 2217.9 & \textbf{0.4332} & 0.5546 \\
8 & 2{,}465{,}379 & 0.0759 & 2.611 & 1975.2 & 0.4266 & 0.5516 \\
\bottomrule
\end{tabular}
\end{fitwidth}
\end{minipage}
\end{center}

\begin{figure}[H]
\centering
\includegraphics[width=\linewidth]{figures/kmeans_elbow_silhouette.png}
\caption{Left: Elbow curve ($k$ vs WCSS). Right: silhouette score vs $k$.}
\label{fig:elbow}
\end{figure}

\textbf{Selecting and justifying $k$.} The two criteria disagree, and the
disagreement is itself the finding:
\begin{itemize}[nosep]
\item \textbf{Silhouette and the elbow both point to a small $k$.} The
silhouette score is maximised at $k = 2$ (0.3946) and falls monotonically
afterwards; the sharpest bend in the WCSS curve (largest second difference)
is at $k = 3$. Every internal metric agrees: Davies--Bouldin is lowest
(best) and Calinski--Harabasz highest (best) at $k = 2$.
\item \textbf{External agreement peaks near $k = 6$--$7$.} ARI is maximised
at $k = 7$ (0.4332) with $k = 6$ almost identical (0.4196), and NMI peaks
exactly at $k = 6$ (0.5593).
\end{itemize}
\noindent \textbf{$k = 2$ is the honest unsupervised answer} -- it is what the
geometry of the data actually supports -- while \textbf{$k = 6$ is the answer
that matches the labelling}. Both solutions are therefore carried forward.

\begin{center}
\begin{minipage}{\linewidth}\centering
\textbf{Table 2: What the $k = 2$ Solution Actually Found}\par\smallskip
\begin{fitwidth}
\begin{tabular}{lcc}
\toprule
True activity & Cluster 0 & Cluster 1 \\
\midrule
LAYING             & 1932 & 12 \\
SITTING            & 1774 & 3 \\
STANDING           & 1906 & 0 \\
WALKING            & 0 & 1722 \\
WALKING\_DOWNSTAIRS & 0 & 1406 \\
WALKING\_UPSTAIRS   & 8 & 1536 \\
\bottomrule
\end{tabular}
\end{fitwidth}
\end{minipage}
\end{center}
\textbf{Inference:} The $k = 2$ solution is an almost perfect static-vs-dynamic
separation: only 23 of 10{,}299 windows (0.22\%) fall on the wrong side. This
is why its silhouette is so much higher than any other $k$ -- it is cutting
along the one genuinely wide gap in the feature space. Its ARI is only 0.3296
purely because it is being scored against a 6-class ground truth that it was
never trying to reproduce.

\section{Model B: DBSCAN}
\label{sec:dbscan}

\textbf{Why DBSCAN is run on a PCA-reduced space.} DBSCAN was first run on the
full 561-dimensional standardized matrix. The results demonstrate the curse of
dimensionality directly: at $\epsilon = 5.0$ and $\epsilon = 7.5$ \emph{every
single point} is labelled noise (no point has 10 neighbours within the
radius), and even at $\epsilon = 12.5$ the best result still discards 44.5\%
of the data as noise with ARI 0.2207. In 561 dimensions all pairwise distances
concentrate into a narrow band, so no single global radius can distinguish
dense regions from sparse ones. The tuning grid was therefore repeated on the
10-component PCA projection.

\begin{center}
\begin{minipage}{\linewidth}\centering
\textbf{Table 3: DBSCAN Tuning on the 10-D PCA Space (selected rows)}\par\smallskip
\begin{fitwidth}
\begin{tabular}{cccccc}
\toprule
$\epsilon$ & \textit{minPts} & Clusters & Noise Fraction & ARI & NMI \\
\midrule
1.5 & 5 & 25 & 88.1\% & 0.0088 & 0.1297 \\
2.5 & 5 & 53 & 52.6\% & 0.1857 & 0.3309 \\
3.0 & 5 & 45 & 38.6\% & 0.2399 & 0.3880 \\
4.0 & 5 & 30 & 18.1\% & 0.3082 & 0.4558 \\
5.0 & 5 & 13 & 7.2\% & 0.3221 & 0.4992 \\
\textbf{6.0} & \textbf{5} & \textbf{7} & \textbf{3.2\%} & \textbf{0.3274} & \textbf{0.5245} \\
6.0 & 10 & 4 & 5.0\% & 0.3249 & 0.5150 \\
6.0 & 20 & 3 & 8.1\% & 0.3210 & 0.4996 \\
8.0 & 5 & 1 & 1.1\% & 0.0013 & 0.0118 \\
\bottomrule
\end{tabular}
\end{fitwidth}
\end{minipage}
\end{center}

\begin{figure}[H]
\centering
\includegraphics[width=0.8\linewidth]{figures/dbscan_tuning.png}
\caption{DBSCAN tuning: Adjusted Rand Index vs $\epsilon$ for three
\textit{minPts} settings on the 10-D PCA space.}
\label{fig:dbscan}
\end{figure}
\textbf{Inference:} DBSCAN's behaviour is dominated by
$\epsilon$, and the sweep brackets the optimum cleanly. Below
$\epsilon = 2.5$ the radius is smaller than the typical nearest-neighbour
distance, so 50--100\% of the data is discarded as noise and the few surviving
clusters are fragments. Raising $\epsilon$ steadily merges those fragments and
recovers structure, peaking at $\epsilon = 6.0$, \textit{minPts} $= 5$, which
labels only 3.2\% of windows as noise and reaches ARI 0.3274 / NMI 0.5245 with
7 clusters -- close to the 6 real activities. Past that point the radius is
wide enough to bridge the static/dynamic gap and the whole dataset collapses
into a \emph{single} cluster at $\epsilon = 8.0$ (ARI 0.0013). The
\textit{minPts} parameter acts as a smoother rather than a driver: at fixed
$\epsilon = 6.0$, raising it from 5 to 20 merely merges the smaller clusters
(7 $\rightarrow$ 3) and raises the noise fraction from 3.2\% to 8.1\%, with
almost no change in ARI.

\section{Model C: Hierarchical Agglomerative Clustering}

Ward linkage requires the full pairwise distance structure, so the linkage
comparison and the dendrogram are computed on a stratified subsample of
3{,}000 windows (500 per activity). The same subsample is used for every
linkage, so the rows of Table~4 are directly comparable.

\begin{center}
\begin{minipage}{\linewidth}\centering
\textbf{Table 4: Linkage Criterion Comparison (HAC, $k = 6$, 3{,}000-window subsample)}\par\smallskip
\begin{fitwidth}
\begin{tabular}{lccccc}
\toprule
Linkage & Silhouette & Davies--Bouldin & Calinski--Harabasz & ARI & NMI \\
\midrule
\textbf{Ward} & 0.0809 & 2.726 & \textbf{697.6} & \textbf{0.3434} & \textbf{0.4946} \\
Complete & 0.3851 & 1.069 & 63.1 & 0.0009 & 0.0290 \\
Average & 0.4193 & 1.155 & 63.6 & 0.0009 & 0.0294 \\
Single & \textbf{0.5542} & \textbf{0.267} & 17.4 & 0.0000 & 0.0033 \\
\bottomrule
\end{tabular}
\end{fitwidth}
\end{minipage}
\end{center}
\textbf{Inference:} This table is the clearest cautionary tale in the
experiment: \textbf{the linkage with the best internal scores is the worst
clustering}. Single linkage achieves a silhouette of 0.5542 and a
Davies--Bouldin of 0.267 -- numbers that look excellent in isolation -- while
scoring an ARI of essentially zero. It does this by the classic single-linkage
failure mode: chaining almost every point into one giant cluster and splitting
off a handful of outliers as the remaining five. Complete and average linkage
fail the same way, less extremely. Only Ward linkage, which explicitly
minimises within-cluster variance and therefore resists chaining, produces a
clustering that corresponds to the real activities (ARI 0.3434), and it does
so while posting the \emph{lowest} silhouette of the four.

\begin{figure}[H]
\centering
\includegraphics[width=\linewidth]{figures/dendrogram_ward.png}
\caption{Dendrogram (Ward linkage, last 30 merges, 3{,}000-window subsample).}
\label{fig:dendro}
\end{figure}
\textbf{Inference:} The dendrogram's two highest merges are separated from
everything below them by a large vertical gap -- the same static/dynamic split
K-Means found at $k = 2$. Cutting below that gap yields the 6-cluster solution
used in Table~5.

\section{Comparison of All Three Algorithms}

\begin{center}
\begin{minipage}{\linewidth}\centering
\textbf{Table 5: Clustering Algorithm Comparison}\par\smallskip
\begin{fitwidth}
\begin{tabular}{lccccccc}
\toprule
Algorithm & Clusters & Noise & Silhouette & Davies--Bouldin & Calinski--Harabasz & ARI & NMI \\
\midrule
K-Means ($k=2$) & 2 & 0\% & 0.3946 & 1.071 & 7880.8 & 0.3296 & 0.5455 \\
K-Means ($k=6$) & 6 & 0\% & 0.1095 & 2.384 & 2556.5 & \textbf{0.4196} & \textbf{0.5593} \\
DBSCAN ($\epsilon=6.0$, \textit{minPts}$=5$) & 7 & 3.2\% & \textbf{0.4151} & \textbf{0.770} & 3449.5 & 0.3274 & 0.5245 \\
Hierarchical (Ward, $k=6$) & 6 & 0\% & 0.0809 & 2.726 & 697.6 & 0.3434 & 0.4946 \\
\bottomrule
\end{tabular}
\end{fitwidth}
\end{minipage}
\end{center}

\begin{figure}[H]
\centering
\includegraphics[width=\linewidth]{figures/algorithm_comparison.png}
\caption{External validation (ARI, NMI) and internal validation (silhouette)
for all four clustering solutions.}
\label{fig:comparison}
\end{figure}

\begin{figure}[H]
\centering
\includegraphics[width=\linewidth]{figures/cluster_visualisation.png}
\caption{Cluster assignments in 2-D. Top row: PCA projection. Bottom row:
t-SNE projection. Left to right: true activities, K-Means, DBSCAN.}
\label{fig:clusters}
\end{figure}

\begin{figure}[H]
\centering
\includegraphics[width=\linewidth]{figures/cluster_activity_contingency.png}
\caption{Cluster-to-activity contingency tables for the three algorithms.}
\label{fig:contingency}
\end{figure}

\section{Observation Questions}
\begin{itemize}
\item \textbf{Which algorithm produced the most meaningful clusters? Why?}
It depends on what "meaningful" is measured
against, and the two answers are genuinely different.
\textbf{Against the activity labels, K-Means with $k = 6$ wins} (ARI 0.4196,
NMI 0.5593) -- no other configuration matches the labelling as closely.
\textbf{Against the geometry of the data, the 2-cluster solutions win}: the
$k = 2$ K-Means split separates static from dynamic activities with 99.8\%
purity and posts the best Calinski--Harabasz score by a factor of two.
DBSCAN at $\epsilon = 6.0$ is arguably the best all-round compromise: it finds
7 clusters without being told how many, keeps 96.8\% of the data, and achieves
both the best silhouette (0.4151) and the best Davies--Bouldin (0.770) of any
solution while staying competitive on NMI (0.5245). The honest summary is that
all three algorithms find the same dominant structure -- motion versus
stillness -- and all three struggle with the same boundary, SITTING versus
STANDING, which is nearly invisible to accelerometer features because both are
motionless postures differing only in orientation.
\item \textbf{How sensitive was K-Means to the choice of $k$?} Very sensitive, and in opposite directions
depending on the metric. Internal quality collapses as $k$ grows: the
silhouette falls from 0.3946 at $k = 2$ to 0.0759 at $k = 8$, a 5-fold
decrease, because each additional split has to cut through a genuinely dense
region. External agreement does the opposite, and non-monotonically: ARI dips
from 0.3296 ($k=2$) to 0.2912 ($k=5$), then jumps sharply to 0.4196 at
$k = 6$. That jump is the signature of $k = 6$ finally having enough clusters
to place one in each activity region. So the choice of $k$ does not merely
change cluster granularity, it changes which structure is recovered -- and the
elbow and silhouette heuristics alone would have led to $k = 2$ or $k = 3$,
neither of which recovers the activities.
\item \textbf{Did DBSCAN detect noise or small clusters effectively?} Yes, once given a space it could work in. The
noise fraction is a direct, tunable dial: 88.1\% at $\epsilon = 1.5$ falling to
3.2\% at $\epsilon = 6.0$. At the chosen setting DBSCAN isolates 333 windows
(3.2\%) as noise while keeping 7 clusters, and those clusters earn the highest
silhouette of any method here -- evidence that excluding the sparse boundary
points genuinely sharpens the remaining structure, which is exactly what
DBSCAN is for and what neither K-Means nor HAC can do (both are forced to
assign every point). It also found small clusters that the partitional methods
cannot express: the $\epsilon = 5.0$ solution splits into 13 clusters, several
of them small sub-groups within the dynamic activities. The important
qualification is that \textbf{DBSCAN completely failed on the raw
561-dimensional data} -- 100\% noise at $\epsilon \le 7.5$, and 44.5\% noise at
best -- so its success is conditional on the PCA reduction, not intrinsic.
\item \textbf{How does linkage choice (single/complete/ward) affect
hierarchical clustering?} Dramatically -- it is the difference between a
working clustering and a useless one. Ward linkage achieves ARI 0.3434, while
complete (0.0009), average (0.0009) and single (0.0000) linkage all produce
essentially random labellings. Single, complete and average linkage merge on
extremal point-to-point distances, so on a dataset with a dense core and a
sparse fringe they chain nearly everything into one cluster and split off tiny
outlier groups. Ward instead merges whichever pair minimises the increase in
total within-cluster variance, a criterion that resists chaining and prefers
balanced, compact clusters -- which matches how these activities actually sit
in feature space.
\item \textbf{Which internal metric best matched your visual intuition of
cluster quality?} None of them matched it reliably, and that is
the most useful lesson in this experiment. \textbf{Silhouette was actively
misleading in the linkage comparison}: single linkage scored 0.5542, the best
of all four linkages, for a clustering that is meaningless (ARI 0.0000), while
Ward scored the worst (0.0809) for the only usable one. Davies--Bouldin agreed
with silhouette and was equally wrong (0.267 for single linkage). The one
internal metric that tracked visual intuition most closely was
\textbf{Calinski--Harabasz}, which ranked Ward (697.6) an order of magnitude
above single linkage (17.4) and correctly preferred the visually obvious
2-cluster split (7880.8). The general pattern is that silhouette and
Davies--Bouldin reward compact, well-separated clusters \emph{whatever those
clusters contain} -- so one giant cluster plus a few tight outlier groups can
score beautifully. Internal metrics are best read alongside a projection plot
and, where labels exist, alongside ARI/NMI.
\end{itemize}

\section{Conclusion}
Three clustering algorithms were applied to 10,299 windows of
smartphone sensor data described by 561 engineered features, with the six
activity labels held back and used only for external validation.

All three recover the same dominant structure, and it is not the six-activity
structure the labels describe: the data's strongest signal is a two-way split
between static postures (LAYING, SITTING, STANDING) and dynamic gaits
(WALKING, WALKING\_UPSTAIRS, WALKING\_DOWNSTAIRS). K-Means finds that split at
$k = 2$ with 99.8\% purity, the Ward dendrogram shows it as the highest and
widest merge, and DBSCAN's collapse into one cluster at $\epsilon = 8.0$ marks
the exact radius at which the gap is bridged. Measured against the labels,
K-Means at $k = 6$ is the best performer (ARI 0.4196, NMI 0.5593); measured on
internal quality, DBSCAN at $\epsilon = 6.0$ is (silhouette 0.4151,
Davies--Bouldin 0.770, with only 3.2\% of points discarded as noise).

The residual errors are consistent across every method and have a physical
explanation: the $k = 6$ contingency table shows SITTING and STANDING merged
into one cluster (1,238 and 1,346 windows in cluster 0), because both are
motionless and accelerometer-derived features cannot separate two stationary
postures that differ only in trunk orientation. The walking variants, by
contrast, are separated reasonably well.

Two methodological findings are worth carrying forward. First, DBSCAN is
unusable on the raw 561-dimensional space -- every point is labelled noise at
small radii because distances concentrate in high dimensions -- and only
becomes competitive after PCA reduction to 10 components. Second, internal
validation metrics can be actively misleading: single linkage posted the best
silhouette and Davies--Bouldin scores in the entire experiment while producing
a clustering with an ARI of zero. Cluster quality should be judged from
projections and, when labels exist, external metrics -- not from a single
internal score.

\section{References}
\begin{enumerate}[nosep]
\item D. Anguita, A. Ghio, L. Oneto, X. Parra, and J. L. Reyes-Ortiz, ``A
Public Domain Dataset for Human Activity Recognition Using Smartphones,''
\emph{21st European Symposium on Artificial Neural Networks (ESANN)}, 2013.
\item M. Ester, H.-P. Kriegel, J. Sander, and X. Xu, ``A Density-Based
Algorithm for Discovering Clusters in Large Spatial Databases with Noise,''
\emph{KDD}, pp. 226--231, 1996.
\item J. H. Ward Jr., ``Hierarchical Grouping to Optimize an Objective
Function,'' \emph{Journal of the American Statistical Association}, vol. 58,
no. 301, pp. 236--244, 1963.
\item P. J. Rousseeuw, ``Silhouettes: A graphical aid to the interpretation and
validation of cluster analysis,'' \emph{Journal of Computational and Applied
Mathematics}, vol. 20, pp. 53--65, 1987.
\item Scikit-learn Developers, ``Clustering,'' scikit-learn documentation.
[Online]. Available: \url{https://scikit-learn.org/stable/modules/clustering.html}
\end{enumerate}

\end{document}
